%! Function Operations — Unit 1, Lesson 1.2 — AP Practice (ANSWER KEY)
% EXTRA CREDIT — optional, exactly TWO pages, placed between the notes and the graded homework.
% Page 1 is Section I, five multiple-choice items with four options (A)–(D); page 2 is
% Section II, one multi-part free-response question on a table. Its contexts (a school concert,
% a bike shop) appear nowhere else in the lesson.
% Distractors are the lesson's errors on purpose: MC 1 adds or reverses when a difference was
% asked, MC 2 answers (g-f) when (f-g) was asked, MC 3 reads the domain off the SIMPLIFIED
% quotient (the lesson's crux), MC 4 adds when a product was asked (and drops a middle term),
% MC 5 forgets that a parent's own excluded input is excluded from the sum as well.
%
% PAGE LOCKSTEP with the other half of the pair: the key marks the correct option in its
% LABEL (no text added to the line) and prose answers are \writespace{H}{…}, fixed height both
% ways. The explicit \newpage fixes the page break in both files.
\documentclass[10pt]{article}
\usepackage{precalculus-article}
\usepackage{precalculus-key}
\newcommand{\TallMath}[1]{$\displaystyle #1\rule[-1.4em]{0pt}{3.2em}$}

\begin{document}

\pageheader{Unit 1, Lesson 1.2}{AP Practice --- Answer Key}

\vspace{0.12in}

\boxguard[8]
\begin{remindbox}
\small
\textbf{Extra credit --- optional.} These two pages are written in the style of the AP
Precalculus exam. Your graded homework is the last section of this packet; do it first. Every
answer choice below is a mistake someone really makes, so read all four before you choose.
\end{remindbox}

\vspace{0.12in}

\begin{headlinebox}{goldbg}
\small
\textbf{Section I --- Multiple Choice.} Circle the letter of the single best answer. No calculator.
\end{headlinebox}

\vspace{0.06in}

\begin{enumerate}[leftmargin=1.9em, itemsep=12pt]

\item A school concert sells $n$ tickets. The revenue is $T(n) = 24n$ dollars and the cost of
      staging the concert is $C(n) = 8n + 450$ dollars. Which function gives the concert's
      profit for $n$ tickets?
      \begin{enumerate}[leftmargin=2.6em, itemsep=2pt]
        \item[(A)] $(T + C)(n) = 32n + 450$
        \item[{\color{keyred}$\checkmark$\,(B)}] $(T - C)(n) = 16n - 450$
        \item[(C)] $(C - T)(n) = 450 - 16n$
        \item[(D)] $(TC)(n) = 192n^2 + 10800n$
      \end{enumerate}

\item \begin{minipage}[t]{0.52\linewidth}
      The table gives values of the functions $f$ and $g$. What is the value of $(f - g)(2)$?
      \begin{enumerate}[leftmargin=2.6em, itemsep=2pt]
        \item[(A)] $-4$
        \item[(B)] $10$
        \item[(C)] $21$
        \item[{\color{keyred}$\checkmark$\,(D)}] $4$
      \end{enumerate}
      \end{minipage}\hfill
      \begin{minipage}[t]{0.42\linewidth}
      \vspace{-2pt}
      \centering
      {\renewcommand{\arraystretch}{1.3}
      \begin{tabular}{|c|c|c|c|}
      \hline
      $x$    & 1 & 2 & 3 \\ \hline
      $f(x)$ & 5 & 7 & 2 \\ \hline
      $g(x)$ & 6 & 3 & 9 \\ \hline
      \end{tabular}}
      \end{minipage}

\item Let $f(x) = x^2 - 49$ and $g(x) = x - 7$. What is the domain of $\dfrac{f}{g}$?
      \begin{enumerate}[leftmargin=2.6em, itemsep=2pt]
        \item[(A)] All real numbers
        \item[{\color{keyred}$\checkmark$\,(B)}] All real numbers except $7$
        \item[(C)] All real numbers except $-7$
        \item[(D)] All real numbers except $7$ and $-7$
      \end{enumerate}

\item If $f(x) = x + 3$ and $g(x) = x - 2$, which of the following is $(fg)(x)$?

      \smallskip
      \textbf{(A)} $2x + 1$ \hspace{2.0em} \textbf{(B)} $x^2 - 6$ \hspace{2.0em} {\color{keyred}\textbf{$\checkmark$\,(C)}} $x^2 + x - 6$ \hspace{2.0em} \textbf{(D)} $x^2 - x - 6$

\item Let $f(x) = \dfrac{1}{x - 4}$ and $g(x) = x^2$. What is the domain of $f + g$?
      \begin{enumerate}[leftmargin=2.6em, itemsep=2pt]
        \item[{\color{keyred}$\checkmark$\,(A)}] All real numbers except $4$
        \item[(B)] All real numbers
        \item[(C)] All real numbers except $0$
        \item[(D)] All real numbers except $4$ and $0$
      \end{enumerate}

\end{enumerate}

\newpage

\begin{headlinebox}{goldbg}
\small
\textbf{Section II --- Free Response.} Show your reasoning. Answers about the bike shop must be
written \emph{in context}, with units --- that is what earns credit on the AP exam.
\end{headlinebox}

\vspace{0.08in}

\begin{enumerate}[leftmargin=1.9em, itemsep=10pt, start=6]

\item A bike shop rents bikes by the day. For $x$ bikes rented on a given day, $R(x)$ is the
      shop's revenue in dollars and $C(x)$ is its operating cost in dollars.

      \smallskip
      \centerline{%
      {\renewcommand{\arraystretch}{1.3}
      \begin{tabular}{|c|c|c|c|c|c|}
      \hline
      $x$ (bikes) & 0 & 5 & 10 & 15 & 20 \\ \hline
      $R(x)$ (\$) & 0 & 90 & 180 & 270 & 360 \\ \hline
      $C(x)$ (\$) & 120 & 135 & 150 & 165 & 180 \\ \hline
      \end{tabular}}}

      \smallskip
      \textbf{(a)} Find $(R - C)(10)$ and interpret it in the context of the problem, with units.
      \writespace{2.0cm}{$(R-C)(10) = 180 - 150 = 30$. On a day when the shop rents 10 bikes, it takes in \$30 more than it spends --- a profit of 30 dollars.}

      \textbf{(b)} Find $(C - R)(5)$. Explain how this value is related to $(R - C)(5)$, and why
      the two are not interchangeable.
      \writespace{2.2cm}{$(C-R)(5) = 135 - 90 = 45$, while $(R-C)(5) = -45$. They are opposites. $C-R$ is the money the shop is short; $R-C$ is its profit, which is negative here. Reversing the order reverses the meaning.}

      \textbf{(c)} The shop's profit function is $P = R - C$. Using the table, between which two
      consecutive $x$-values in the table does $P$ first become positive? Justify your answer
      with values.
      \writespace{2.4cm}{$P(0) = -120$, $P(5) = -45$, $P(10) = 30$, $P(15) = 105$, $P(20) = 180$. $P$ first becomes positive between $x = 5$ and $x = 10$: it is still negative at 5 bikes and positive at 10.}

      \textbf{(d)} Let $d(x) = x$, so that the average revenue per bike is the quotient
      $\left(\dfrac{R}{d}\right)(x)$. Explain why this quotient has no value at $x = 0$, and
      state its domain for the $x$-values in the table.
      \writespace{2.4cm}{At $x = 0$ the bottom function is $d(0) = 0$, and division by zero is undefined --- with no bikes rented there is no ``per bike'' amount. The domain is the table's values except 0: $x = 5, 10, 15, 20$.}

      \textbf{(e)} The shop's rates change so that $R(x) = 18x$ and $C(x) = 3x + 120$ for
      $0 \le x \le 20$ bikes. Write $P(x) = (R - C)(x)$ in simplified form, find the smallest
      whole number of bikes for which $P(x) > 0$, and state the domain of $P$ in this context.
      \writespace{2.4cm}{$P(x) = 18x - (3x + 120) = 15x - 120$. $P(x) > 0$ when $15x > 120$, that is $x > 8$, so the smallest whole number is 9 bikes. Domain: whole numbers of bikes with $0 \le x \le 20$.}

\end{enumerate}

\end{document}
